\section{The single-band completion barrier}
\label{sec:single-band-barrier}

We now compare the literal global prime band inherited by TPC--47 with the
critical very-high face of TPC--59.  The comparison is made at the same
fixed shift and on the same target interval.  It has two parts.  First, an
exact constant-factor sandwich identifies which cofactors one band can
see.  Second, the TPC--61 harmonic certificate is evaluated on every full
dyadic block that the band omits.  Neither part estimates the actual mass
of those blocks.

\subsection{An exact constant-factor window}

Let every target satisfy
\begin{equation}
 N_-\le n=mj+h_0=pr\le N_+,
 \qquad
 \kappa_N:=\frac{N_+}{N_-}.
 \label{eq:band-target-endpoints}
\end{equation}
The fixed support ratios give \(1\le\kappa_N\ll1\).  For a parameter
\(R_0\) in the admissible TPC--45 range, the inherited global prime band is
\begin{equation}
 \cP_{R_0}^{\sharp}
 :=\left\{p\in\Pp:
      \sqrt{N_+}<p\le\frac{N_-}{R_0}\right\}.
 \label{eq:band-single-prime-band}
\end{equation}
Let \(y\) be very-high admissible.  In particular,
\(y>\sqrt{N_+}\), and put
\begin{equation}
 R_{\rm safe}(y):=\frac{N_-}{2y},
 \qquad
 R_{\max}(y):=\frac{N_+}{y}
 =2\kappa_NR_{\rm safe}(y).
 \label{eq:band-safe-max}
\end{equation}
These are the exact TPC--59 safe and maximal cofactor scales
\citep{WangTPC59}.

Define the portion of the native high carrier selected solely by the prime
band condition:
\begin{equation}
 \cW_{R_0}^{\rm band}(y)
 :=\{w\in\cW_H(y):p(w)\in\cP_{R_0}^{\sharp}\}.
 \label{eq:band-high-intersection}
\end{equation}
All other row, orbit, opposite-row, mask, profile, and terminal labels are
those already present in \(\cW_H(y)\); no favorable support is added.

\begin{theorem}[Exact single-band/high-face sandwich]
\label{thm:band-exact-window-sandwich}
For every very-high admissible \(y\),
\begin{equation}
 \boxed{
 \begin{aligned}
 \{w\in\cW_H(y):\ &\kappa_NR_0\le r(w)<R_{\max}(y)\}
 \\
 &\subseteq \cW_{R_0}^{\rm band}(y)
 \\
 &\subseteq
 \{w\in\cW_H(y):R_0\le r(w)<R_{\max}(y)\}.
 \end{aligned}}
 \label{eq:band-exact-window-sandwich}
\end{equation}
Thus the only uncertainty created by the target-width constants lies in
the fixed-ratio transition window
\begin{equation}
 R_0\le r<\kappa_NR_0.
 \label{eq:band-transition-window}
\end{equation}
In particular, every high atom with \(r<R_0\) is absent from the single
band, while every high atom with \(r\ge\kappa_NR_0\) is present in it.
\end{theorem}

\begin{proof}
If \(w\in\cW_{R_0}^{\rm band}(y)\), then
\(p(w)\le N_-/R_0\) and \(n(w)\ge N_-\), so
\begin{equation}
 r(w)=\frac{n(w)}{p(w)}\ge R_0.
 \label{eq:band-right-inclusion-lower}
\end{equation}
Also \(p(w)>y\) and \(n(w)\le N_+\), hence
\begin{equation}
 r(w)<\frac{N_+}{y}=R_{\max}(y).
 \label{eq:band-right-inclusion-upper}
\end{equation}
This proves the right inclusion.

Conversely, let \(w\in\cW_H(y)\) and suppose
\(r(w)\ge\kappa_NR_0\).  Then
\begin{equation}
 p(w)=\frac{n(w)}{r(w)}
 \le\frac{N_+}{\kappa_NR_0}
 =\frac{N_-}{R_0}.
 \label{eq:band-left-inclusion-upper}
\end{equation}
The very-high inequality \(p(w)>y>\sqrt{N_+}\) supplies the lower endpoint
of \eqref{eq:band-single-prime-band}.  Thus \(p(w)\) lies in the band.
The upper bound \(r(w)<R_{\max}(y)\) is automatic on the high face.
\end{proof}

\begin{corollary}[Empty, top-window, and multiscale regimes]
\label{cor:band-three-regimes}
The sandwich has the following exact consequences.
\begin{enumerate}[label=\textup{(\roman*)},leftmargin=3em]
 \item If \(R_0\ge R_{\max}(y)\), then
       \(\cW_{R_0}^{\rm band}(y)=\varnothing\).
 \item If \(R_0\asymp R_{\max}(y)\), the intersection is confined to
       \(O(1)\) top dyadic cofactor blocks.
 \item If \(R_0=X^{\theta_0+o(1)}\) and
       \(R_{\max}(y)=X^{\theta_J+o(1)}\) with
       \(0\le\theta_0<\theta_J\), then the band may contain high blocks
       of exponent \(\theta\ge\theta_0\), but it omits every full block
       with \(2R<R_0\).
\end{enumerate}
These are support statements.  In cases \textup{(ii)} and \textup{(iii)},
the permitted blocks may still carry zero actual atomic mass.
\end{corollary}

\subsection{The endpoint mismatch with the inherited freezing certificate}

At the critical cutoff of TPC--59,
\begin{equation}
 y=C_{\rm vh}Q,
 \qquad
 R_{\max}(y)\asymp J
 =X^{\theta_J+o(1)},
 \qquad
 \theta_J:=\frac{133}{400}.
 \label{eq:band-critical-ceiling}
\end{equation}
The balanced global band of TPC--45 has
\begin{equation}
 R_0=\ell_+
 =X^{\lambda_L+o(1)},
 \qquad
 \lambda_L:=\frac{99979}{210000},
 \label{eq:band-balanced-R0}
\end{equation}
and therefore
\begin{equation}
 \lambda_L-\theta_J
 =\frac{15077}{105000}>0.
 \label{eq:band-balanced-critical-gap}
\end{equation}

\begin{corollary}[The balanced band misses the critical high face]
\label{cor:band-balanced-empty}
For all sufficiently large \(X\), the balanced band satisfies
\begin{equation}
 R_0>R_{\max}(C_{\rm vh}Q),
 \qquad
 \boxed{\cW_{R_0}^{\rm band}(C_{\rm vh}Q)=\varnothing.}
 \label{eq:band-balanced-empty}
\end{equation}
\end{corollary}

This is not peculiar to the exact balanced choice.  Write
\(R_0=X^{\theta_0+o(1)}\).  The sufficient TPC--45 band/cutoff freezing
condition has the form
\begin{equation}
 (\lambda_L-\theta_0)_+
 +\lambda_0\left(1-\frac1A\right)<\frac1{400},
 \qquad A>1,
 \label{eq:band-inherited-freezing-condition}
\end{equation}
where the second term is nonnegative \citep{WangTPC45}.  Hence every band
certified by this inequality obeys
\begin{equation}
 \theta_0>\lambda_L-\frac1{400}
 =\theta_J+\frac{29629}{210000}.
 \label{eq:band-certified-exponent-gap}
\end{equation}

\begin{theorem}[Single-band endpoint incompatibility]
\label{thm:band-endpoint-incompatibility}
At \(y=C_{\rm vh}Q\), every single band certified by
\eqref{eq:band-inherited-freezing-condition} is eventually disjoint from
the critical very-high face.  Conversely, retuning
\(R_0=X^{\theta_J+o(1)}\) to meet the top critical cofactor window incurs
the coefficient-blind band exponent
\begin{equation}
 \lambda_L-\theta_J
 =\frac{15077}{105000}
 >\frac1{400},
 \label{eq:band-retuned-cost}
\end{equation}
so the inherited sufficient freezing certificate no longer closes.
\end{theorem}

\begin{proof}
Equation \eqref{eq:band-certified-exponent-gap} gives
\(R_0/R_{\max}(y)=X^{c+o(1)}\) for a fixed \(c>0\).  Apply
\cref{cor:band-three-regimes}\textup{(i)}.  For the retuned value,
substitution of \(\theta_0=\theta_J\) in the first term of
\eqref{eq:band-inherited-freezing-condition} gives
\eqref{eq:band-retuned-cost}; adding the nonnegative cutoff term cannot
improve the inequality.
\end{proof}

The theorem closes one simultaneous use of the old support-only freezing
certificate and the critical high face.  It does not say that a different
coefficient-specific freezing theorem is impossible.

\subsection{Full omitted blocks and the harmonic zero-saving theorem}

Suppose now that the band is retuned so that
\(R_0<R_{\max}(y)\).  The exact sandwich still forces every cofactor
\(r<R_0\) to be absent from the band-generated relation.  Let
\([R,2R)\) be a full dyadic block satisfying
\begin{equation}
 1\le R<2R\le R_0,
 \qquad
 R\le R_{\max}(y).
 \label{eq:band-full-omitted-block}
\end{equation}
For a row \(m\), call the block \emph{fully omitted} when
\begin{equation}
 [R,2R)\cap\N\subseteq E_m
 :=\{r\le R_{\max}(y):r\notin\cI_m\}.
 \label{eq:band-full-omission}
\end{equation}
The TPC--61 harmonic capacity is
\begin{equation}
 \mathfrak h_m(E_m)
 :=\frac m{\varphi(m)}
   \sum_{r\in E_m}
   \left(1+\frac{J/r}{\log(2+J/r)}\right),
 \label{eq:band-harmonic-capacity}
\end{equation}
and its effective high exposure is
\begin{equation}
 H_m^{\rm eff}=\frac{W_H(m)}{U_m},
 \qquad
 0<H_m^{\rm eff}\le\#\cJ_X\ll J.
 \label{eq:band-effective-exposure-upper}
\end{equation}

\begin{theorem}[Harmonic zero-saving on a full omitted block]
\label{thm:band-harmonic-zero-saving}
Assume \eqref{eq:band-full-omission} on an active row.  Then
\begin{equation}
 \boxed{
 \mathfrak h_m(E_m)
 \gg
 R+\frac{J}{\log(2+J/R)}.}
 \label{eq:band-harmonic-lower}
\end{equation}
Consequently,
\begin{equation}
 \boxed{
 \frac{\mathfrak h_m(E_m)}{H_m^{\rm eff}}
 \gg
 \frac RJ+\frac1{\log(2+J/R)}.}
 \label{eq:band-harmonic-ratio-lower}
\end{equation}
If
\begin{equation}
 J=X^{\theta_J+o(1)},
 \qquad
 R=X^{\theta+o(1)},
 \qquad 0\le\theta\le\theta_J,
 \label{eq:band-block-exponents}
\end{equation}
then the right side of \eqref{eq:band-harmonic-ratio-lower} is
\(X^{-o(1)}\).  Therefore the TPC--61 harmonic sufficient condition cannot
certify
\begin{equation}
 \frac{\mathfrak h_m(E_m)}{H_m^{\rm eff}}
 \le X^{-\sigma+o(1)}
 \label{eq:band-unavailable-harmonic-saving}
\end{equation}
for any fixed \(\sigma>0\).
\end{theorem}

\begin{proof}
The factor \(m/\varphi(m)\) is at least one.  There are \(\gg R\)
integers in \([R,2R)\).  Each contributes one to
\eqref{eq:band-harmonic-capacity}.  Moreover, for
\(R\le r<2R\),
\begin{equation}
 \frac{J/r}{\log(2+J/r)}
 \ge
 \frac{J/(2R)}{\log(2+J/R)}.
 \label{eq:band-harmonic-term-lower}
\end{equation}
Summing over the block proves \eqref{eq:band-harmonic-lower}.  Divide by
\(H_m^{\rm eff}\ll J\) to obtain
\eqref{eq:band-harmonic-ratio-lower}.  Under
\eqref{eq:band-block-exponents}, either \(R/J=X^{-o(1)}\), or
\(\log(2+J/R)=O(\log X)\); in both cases the displayed lower bound is
\(X^{-o(1)}\), which is larger than \(X^{-\sigma+o(1)}\) for every fixed
positive \(\sigma\).
\end{proof}

The same conclusion appears in the exponent law of TPC--61.  A full block
has cardinality exponent \(e(\theta)=\theta\), while necessarily
\(H_m^{\rm eff}\le X^{\theta_J+o(1)}\).  Even at the optimal balanced
exposure exponent \(\eta_H=\theta_J\), its certified saving is
\begin{equation}
 \sigma_{\rm cof}(\theta)
 =\eta_H-e(\theta)-(\theta_J-\theta)
 =0.
 \label{eq:band-zero-saving-exponent}
\end{equation}
Any smaller effective-exposure exponent only worsens the certificate.

\begin{remark}[A certificate barrier, not an actual lower bound]
\label{rem:band-certificate-only}
TPC--61 proves the one-sided implication
\begin{equation}
 \frac{W_M(m)}{W_H(m)}
 \ll\frac{\mathfrak h_m(E_m)}{H_m^{\rm eff}}.
 \label{eq:band-one-sided-harmonic-bound}
\end{equation}
A lower bound for the right side does not lower-bound the left side.  The
omitted block may contain no realized largest-prime cofactors, or its
literal weights may vanish.
\Cref{thm:band-harmonic-zero-saving} says only that the declared
support-cardinality/Brun--Titchmarsh certificate cannot prove a fixed-power
upper saving for a full omitted block.  A weighted largest-prime-factor
distribution theorem, an actual terminal-incidence estimate, or signed
cancellation could still do so.
\end{remark}

The window sandwich and endpoint exponent comparison are L0 arithmetic
support statements specialized at L1 to the literal fixed-\(h_0\) target
interval and prime band.  The harmonic lower calculation closes one L1
support-only certificate.  None of these results proves a positive
high-face mass, a lower terminal square, a M\"obius correlation estimate,
or an L2 fixed-shift parity improvement.
